\cotefile{157-2}
\scanpage{383}\mspage{122}

\[
H_T(Y)_{\mathrm{Cat\,hot}}=\mathrm{Hom}^{\unsure{\deleted{\circ}|\circledast}}_{\mathrm{Ind\,Cat\,h.t}}(Y,H_T)
\]

Pour résumer ces lumières et perplexités sur les $C_\infty(\underline t)$

a) À \unsure{tout} $\underline t=(t_n)_{n\in\mathbb N^{\mathrm o}}\in\underline{\mathrm{Hom\,cof}}(\mathbb N^{\mathrm o},T)$ est associé un système inductif, indexé par $\mathbb N$, de catégories $C_n(\underline t)$ $(=C_{t_n})$, et leur limite inductive

\[
C_\infty(\underline t).
\]

b) Ce\deleted{\unsure{tte construction}}\inserted{système} inductif, et aussi sa limite, est \inserted{foncteur} contravariant en $\underline t$. Si $\underline t,\underline t'\in\mathrm{Ob}\,\underline{\mathrm{Hom\,cof}}$, alors on peut comparer $C_\infty(\underline t)$, $C_\infty(t')$ et les syst. inductifs correspondants, grâce à

\diagram{\begin{tikzcd}[ampersand replacement=\&]
\& \underline t\unsure{*|\times}\underline t' \arrow[dl]\arrow[dr] \& \\
\underline t \& \& \underline t'
\end{tikzcd}}

c) L'objet

\[
\big(C_n(\underline t)\big)_{n\in\mathbb N}\in\mathrm{Ob}\,\mathrm{Ind\,Cat\,hot}
\]

défini par le système inductif $(C_n(\underline t))_{n\in\mathbb N}$ de $\mathrm{Cat}$ est \inserted{(à isom unique près)} indépendant du choix de $\underline t$, soit $H_{\unsure{T|C}}$\inserted{, ou $H$}, Les hom

\[
\underline C_\infty(\underline t')\xrightarrow{u^*}\deleted{\unsure{C}}C_\infty(\underline t)
\tag*{$(\unsure{*|\leftrightarrow})$}
\]

associé à $u:\underline t\longrightarrow\underline t'$ sont compatibles avec cette donnée, i.e. on a comm. \deleted{\unsure{du}}

\diagram{\begin{tikzcd}[ampersand replacement=\&,column sep=small]
\underset{\mathrm{Ind\,Cat\,h.t}}{\underset{\longrightarrow}{\lim}}\;\underline C_*(\underline t') \arrow[rr] \& \& \underset{\mathrm{Ind\,Cat\,h.t}}{\underset{\longrightarrow}{\lim}}\; C_{*}(\underline t)\\
\& H\deleted{\unsure{T}}\arrow[ul,"\sim"']\arrow[ur,"\sim"] \&
\end{tikzcd}}

% archive sheet number: 382
